\documentclass[11pt,confidential]{article}

%\usepackage{axrtechdocs}
\usepackage{axioma}
\usepackage{graphicx}
\usepackage{float}
\usepackage{amsmath}
\usepackage{amsfonts}
\usepackage{amssymb}
\usepackage[capitalise,noabbrev]{cleveref}
\usepackage{booktabs}
%\newcommand{\tabitem}{~~\llap{\textbullet}~~}
%\usepackage{datetime}
\usepackage{longtable}
\usepackage{array}
\usepackage{multirow}
\usepackage{pbox}
%\usepackage{enumitem}
\renewcommand{\arraystretch}{1.3}
\newcommand{\note}[1]{\vspace{0.2cm}\noindent*************************************************************************************** \\ NOTE: {#1} \\ ***************************************************************************************\\ \vspace{0.1cm}}


\title{Implied Volatility Surface Construction}
\author{Manu Maurette - QPRA}
\date{2Q 2019}

\begin{document}
	\maketitle

\newpage

\vspace*{\fill}
\section*{Version Control}

\begin{center}
    \begin{tabular}{ | l | l | l| p{8cm} |}
    \hline
   Version & Date       & Authors       & Comments      \\ \hline
    1.0    & 5/16/2019  & Manu Maurette & First draft.  \\ \hline
		1.1    & 5/28/2019  & Manu Maurette, Emma Li & First version  \\ \hline
           &            &               &               \\ \hline
			     &            &               &               \\ 
					\hline
    \end{tabular}
\end{center}


\pagebreak


\section{Introduction}

The goal of this document is to study the possibility of generating in house volatility surfaces. Currently Axioma uses a Vendor to obtain Volatility Surfaces. These are then used as inputs in AxR pricing models.


\section{Input Data}

We here give a brief description of the necessary inputs for obtaining implied volatilities

\subsection{Option Prices}

For a given security (Stock, Index), for each listed Strike, Expiration and Option Type (Call/Put), the market price is needed. This information in general is provided by data vendors as:
\begin{itemize}
\item \textbf{Best Bid price available}
\item \textbf{Best Ask price available}
\item Last traded date
\item Daily Volume
\item Open interest
\end{itemize}
 
The first two are essential. The last three, instead, \textbf{could be used for cleaning the data and avoiding using illiquid products that might affect the calculations}.

Currently AXIOMA obtains that data from Thomson Reuters in \textit{EJV$\_$derivs} database (available in Pandora, for example)

We will follow the following assumptions:
\begin{itemize}
\item Market price of an option to be $\frac{Bid+Ask}{2}$ (industry standard)
\item No lag for settlement needs to be taken into account. 
\item No filling for missing quotes of option prices.
\end{itemize}


\subsection{Risk Free rates}

Pricing models also assume the existence of a risk-free rate used for discounting. We will assume this curve is the Zero Coupon USD (or the corresponding currency) Swap Rate curve.

Currently, AXIOMA obtains these rates from Thomson Reuters in \textit{MarketData} database (available in Pandora, for example). See Testing section for an example.

In addition, the rates obtained from the database have to be interpolated to the options times to maturities. This will be continued in the Analytics section


\subsection{Dividend Yield}

\subsubsection{Equity}

Currently, AXIOMA obtains these rates from the \textit{Oracle Equity} database. See Testing section for an example. 

\subsubsection{Indexes}

 
In the case of European, Canadian and Asian Indexes, an industry standard is to estimate the dividend as $q = \frac{PV(div)}{T*S}$. 


\textbf{We are still in review on alternative ways to obtain these }


\section{Modeling}

There are several pieces in the process that require modeling decision:

\subsection{Interest Rate curve interpolation}

As mentioned in the previous section, rates have to be interpolated to get rates for the listed options times to maturities.
This can be done with several methods as:

\begin{itemize}
\item Piecewise Constant
\item Piecewise Linear
\item Cubic Spline
\end{itemize}


\subsection{Dividend Yield}

Estimation of the dividend yield also require a modeling decision.

There are several methods to estimate an Index Dividend Yield. For the US case, an industry standard is to use a Call-Put parity based regression. For dividend paying stock, a call put parity relationship holds:

\[S_0 - PV(Div) + P = C + K e^{-rT},\]

\noindent so, by approximating $PV(Div) = S_0(1-e^{-dT})$

\[C(t) -P(t) = S_0 e^{-dT} - K e^{-rT}.\]

By first order approximation of $e^{-rT}=1-rT$:

\[C(t) -P(t) = S_0 (1-dT) - K (1-rT),\]

regrouping:

\[C(t) -P(t) = S_0 - d S_0 T - K- rKT.\]

The following regression is done:

\[C -P = b_0 + b_1 S + b_2 ST + b_3 K + b_4 KT + b_5 D_{BA}.\]

The regression is difference between the price of a call option and the price of a put
option with the same expiration and strike. When calculating this difference, the bid
price of the call is used with the offer price of the put, and vice versa. $D_{BA}$ is a dummy 
variable set equal to 1 if the call option's bid price is used. $S$ is the underlying security's
(index) closing price, $K$ is the strike price of the call and put options, and $T$ is the time
to expiration in years. The regression is calculated using three months of option data
across all strikes and expirations with an exception of contracts expiring in less than 15
days, for a single underlying. According to the principle of put-call parity, the dividend
yield on the underlying index will be approximately equal to the negative of the estimated
parameter $b_2$.






\subsection{Option Analytics}

\subsubsection{American}

Most US stock options have American style exercise. This requires the need of an American Option model implementation. Some common choices are:

\begin{itemize}

\item CRR binomial tree model
\item Modification of binomial tree models
\item LSMC method
\item Axioma pricers (may include the previous)

\end{itemize}

\subsubsection{European}

For European style options (most US indices are of European nature) the natural choice is the Black Scholes model.


\subsection{Root Finding algorithm for obtaining implied volatilities}

Again, this step requires a modeling decision:

\begin{itemize}
\item Bisection
\item Newton Raphson 
\item Jackel implication (***cite lets be rational***) for European style.
\end{itemize}

\subsection{Volatility Surface interpolation/extrapolation}

Once all the implied volatilities, we need to generate a standard volatility surface. Again, this is an interesting mathematical problem. Some algorithms for doing this include:

\begin{itemize}
\item Piecewise linear
\item Cubic spline
\item SVI-BDK (available in AxR)
\item Tension spline / Thin plate spline
\item Gaussian kernel smoothing (Vega weighted)
\end{itemize}



\subsubsection{Gaussian Kernel Smoothing}

We generate a \textit{standard} volatility surface based on Delta and Time to Maturities. 
We define standard deltas and standard TTMs:
\begin{itemize}
\item TTMs [30, 60, 91, 122, 152, 182, 273, 365, 547, 730] (in days)
\item Deltas $\pm$[20, 25, 30, 35, 40, 45, 50, 55, 60, 65, 70, 75, 80] (in \%)
\item Call-Put
\end{itemize}

For each grid point $j$ (Delta, Time to Maturity, Type), which we perform a smoothing algorithm: 

\[\hat{\sigma_j}=\frac{\sum_i V_i\sigma_i\Phi(x_{ij},y_{ij},z_{ij})}{\sum_i V_i\Phi(x_{ij},y_{ij},z_{ij})}\]

\noindent where $i$ is indexed over all listed options for that day (after the data cleanse); $V_i$ is the option
vega, $\sigma_i$ is the implied volatility and $\Phi()$ is a kernel function:

\[\Phi(x,y,z) = \frac{1}{\sqrt{2\pi}}e^{-\left[\frac{x^2}{2h_1}+\frac{y^2}{2h_2}+\frac{z^2}{2h_3}\right]}.\]
\noindent with:
\begin{itemize}
\item $x_{ij}=\log(T_i/T_j))$
\item $y_{ij} = \Delta_i-\Delta_j$
\item $z_{ij} = I_{\{CP_i=CP_j\}}$
\end{itemize}
The kernel bandwidth parameters are chosen empirically: $h_1=0.05$, $h_1=0.005$, $h_3=0.001$.





\section{Testing and Validation}


\begin{itemize}
\item Replicate implied volatilities for equity index: SPX, DJIA, NASDAQ, FTSE, Russell, Nikkei
\item Implication only
\item Full replication 
\item Replicate implied vols for single-name stock
\item Implication only
\item Full replication 
\end{itemize}

\subsection{First set of tests}

\subsubsection{AAPL stock for $04.04.2019$}
The following describes an end to end tests to generate volatility surface for AAPL, for $04.04.2019$.

\paragraph{EJV Option data}

In SQL, we perform the following Query

\begin{verbatim}
Declare @AsOfDate Date='2019-04-04',
@UndlyRic varchar(100) = 'AAPL.O'
if object_id('tempdb..#EJVDerivsData', 'U') is not null drop table #EJVDerivsData
select bid_px, ask_px, universal_bid_px, universal_ask_px, vol, open_interest, a.ric_root, 
unscaled_strike_px, put_call_indicator, expiration_dt, exercise_style_cd, a.ric
into #EJVDerivsData
from EJV_Derivs.dbo.quote_xref a, EJV_Derivs.dbo.price b 
where underlying_ric=@UndlyRic and put_call_indicator is not null and 
expiration_dt >= @AsOfDate and a.quote_id=b.quote_id and b.trading_dt=@asOfDate
select * from #EJVDerivsData b
\end{verbatim}

\begin{center}
\begin{figure}[h!]
\includegraphics[width=6in]{SQL_EJV}
\caption{EJV derivs query for Option prices}
\end{figure}
\end{center}


We compared option prices (bids and asks) from EJV derivs database and The Vendor and reached to the following conclusion:

\begin{itemize}
\item The same 1604 options were listed in both databases.
\item Both databases have a similar behavior for illiquid options. One (Vendor) puts 0 in the Bid, the alternative puts 0.05 (or 0.01)
\item Removing the above cases (1422 cases), 97\% of the relative Bid differences fall in the (-5\%, 10\%) column of the histogram (see Figure below).

\begin{center}
\begin{figure}[h!]
\includegraphics[width=5in]{BidHistogram_AAPL}
\caption{Bid relative difference histogram (in \%)}
\end{figure}
\end{center}



\item Also, removing the above cases (1422 cases), 89\% of the relative Ask differences fall in the (-8\%, 2\%) 
column of the histogram (see Figure below).


\begin{center}
\begin{figure}[h!]
\includegraphics[width=5in]{AskHistogram_AAPL}
\caption{Ask relative difference histogram (in \%)}
\end{figure}
\end{center}

\end{itemize}


\paragraph{ZC Swap Curve data}

In SQL, we perform the following Query

\begin{verbatim}
SELECT cv.CurrencyEnum, CountryEnum, cv.CurveTypeEnum, CurveShortName, 
CurveLongName, TradeDate, te.Inyears, te.Lud, te.PeriodUnitEnum, te.Description, 
te.Lub, te.NumberOfUnit, cq.Quote
FROM [MarketData].[dbo].[Curve] cv
join [MarketData].[dbo].[CurveNodes] cn on cv.CurveId = cn.CurveId 
join [MarketData].[dbo].[TenorEnum] te on cn.TenorEnum = te.TenorEnum
join [MarketData].[dbo].[CurveNodeQuote] cq on cq.CurveNodeId = cn.CurveNodeId
and TradeDate = '2019-04-04'
and cv.CurrencyEnum ='USD'
and cv.CurveTypeEnum in  ('SwapZC')
order by InYears
\end{verbatim}



\begin{center}
\begin{figure}[h!]
\includegraphics[width=7in]{SQL_ZC}
\caption{MarketData query for SwapZC curve}
\end{figure}
\end{center}


After a cubic spline interpolation of these data, both from the Vendor and the Alternative Axioma sources, we obtain the following plot of the curves:




\begin{center}
\begin{figure}[h!]
\includegraphics[width=7in]{ALTvsVendor_ZCCurve_04042019}
\caption{Differences in ZC Swap curves}
\end{figure}
\end{center}


There is a discrepancy in the first tenors that we currently investigating.



\paragraph{Dividend}

In SQL, we perform the following Query:

\begin{verbatim}
select * from openquery([ORA-GLPROD-marketdb_global], 'select div.dt DivDate, div.value 
from modeldb_global.sub_issue_divyield_active div 
where div.sub_issue_id = ''DQ8L3JXGU111''
and div.dt = TO_DATE(''2019-04-04'', ''YYYY-MM-DD'')') 

\end{verbatim}

\begin{center}
\begin{figure}[h!]
\includegraphics[width=4in]{SQL_Div}
\caption{Oracle Equity query for Dividend yield}
\end{figure}
\end{center}






A comparison was made to how the Vendor computes the dividend. It is not clear how the calculation is done, we can interpret two ways:


\begin{enumerate}
\item $\textnormal{div} =\frac{Divided Payed}{Spot}$ (Past)
\item $\textnormal{div} =\frac{Divided Payed}{Spot}$ (Estimate)
\end{enumerate}

\begin{center}
\begin{figure}[h!]
\includegraphics[width=4in]{DIV_AAPL_04042019}
\caption{Differences in Dividend yield}
\end{figure}
\end{center}











\paragraph{Data cleanse aggregation}

\begin{itemize}

\item Incorporate Options prices 
\item Remove \verb"NAN" values of Options database
\item Generate \verb"TTM" (Time to Maturities) of Expiration dates based on valuation date and Expiration using (ACT/ACT) daycount convention.
\item Generate risk free rate for each TTM from ZC rate data using a \textit{cubic spine interpolation}
\item Incorporate Dividend Yield
\item Incorporate an IsIndex flag
\end{itemize}

Bellow is a picture of part of the DataFrame of cleaned data.

\begin{center}
\begin{figure}[h!]
\includegraphics[width=4in]{CleanDataFrame2}
\caption{Clean data DataFrame}
\end{figure}
\end{center}


\paragraph{Pricing method}

Leisen Reisner tree method was used to price American Option. It is an enhancement of the CRR tree. ***CITE***


\paragraph{IV generation for each listed option in dataframe}

Implied volatilities are generated using a root finding algorithm. Starting from interval $\sigma\in [0,5]$. In this case, \textit{bisection}. Also, the LR algorithm is used in each iteration for pricing.

\paragraph{Generation of model outputs}

At this point, with the obtained Implied Volatility, the following outputs are obtained:

\begin{itemize}
\item  Model Price
\item  Model Delta
\item  Model Vega
\end{itemize}

Here, an extra cleanse is made. Deltas smaller than 5\% and greater than 95\% in abs value are omitted due to potential issues in pricing calculations.
\textbf{We need to study this deeper.}

At this point, it would be possible to do a comparison with the Vendor Database





\paragraph{Generation of volatility surface}

We implemented the Gaussian Kernel Smoothing


\paragraph{Tests}

Comparison tests were performed between the volatility surfaces obtained following these steps, and a volatility surface provided by a vendor.


\begin{center}
\begin{figure}[h!]
\includegraphics[width=7in]{ALTvsVendor_AAPL_04042019}
\caption{Differences in volatilities}
\end{figure}
\end{center}

We see peaks of around 9\% differences for the 20\% Delta (out of the money) and short maturities (1M 3M) and an average of a +-2\% relative difference between both.


\subsubsection{SPX index for $04.04.2019$}

The following describes an end to end tests to generate volatility surface for SPX, for $04.04.2019$.

\paragraph{EJV Option data}

I was provided a Query ***FIND GENERAL QUERY*** with data of listed options.


We compared option prices (bids and asks) from EJV derivs and The Vendor and reached to the following conclusion:

From 18127* options with data both from we found:

\begin{itemize}
\item 1 case with ASKS relative difference 100\%
	They correspond to Asks in Vendor1 = 0

\item 26 (0.14\%) cases with ASKS relative differences above (or equal) 50\%:
	They correspond to Asks in Vendor2 in the set 0.05 and 0.1
	They correspond to Asks in Vendor1 in the set 0.05, 0.1, 0.15 and 24  of them correspond to Bid in Vendor 1 = 0

\item 43 (0.24\%) cases with ASKS relative differences between 25\% and 50\%:
	They correspond to Asks in Vendor2 in the set 0.15, 0.2, 0.25, 0.35 
	They correspond to Asks in Vendor1 in the set 0.1, 0.15, 0.2, 0.25 

\item 183 (1.01\%) cases with ASKS relative differences between 10\% and 25\%:
	All ASKS bellow 3.3 (Vendor1) and 3.7 (Vendor 2)

\item 336 (10.48\%) cases with ASKS relative differences between 5\% and 10\%:
	All ASKS bellow 8.1 (Vendor1) and 8.6 (Vendor 2)

\item 1900 (9.9\%) cases with ASKS relative differences between 1\% and 5\%:

\item 86.27\%  of the cases with ASKS relative differences below 1\%
\end{itemize}



\paragraph{ZC Swap Curve data}

Same as the last case, as it is the same date


\paragraph{Dividend}

We still do not have a clear idea of how to obtain Dividend for Index Options.
We used used $1.86\%$ following several data sources (not from Axioma). 

The vendor uses the regression algorithm described in the Analyitics section.



\paragraph{Data cleanse aggregation}

\begin{itemize}

\item Incorporate Options prices 
\item Generate \verb"TTM" (Time to Maturities) of Expiration dates based on valuation date and Expiration using (ACT/ACT) daycount convention.
\item Generate risk free rate for each TTM from ZC rate data using a \textit{cubic spine interpolation}
\item Incorporate Dividend Yield
\item Incorporate an IsIndex flag (1 in this case)
\end{itemize}



\paragraph{Pricing method}

BSM model was used to price the European style options. 

\paragraph{IV generation for each listed option in dataframe}

Implied volatilities are generated using a root finding algorithm. Starting from interval $\sigma\in [0,8.5]$. In this case, \textit{bisection}. Jeckel algorithm was also tested with the exact same results. 

\paragraph{Generation of model outputs}

At this point, with the obtained Implied Volatility, the following outputs are obtained:

\begin{itemize}
\item  Model Price
\item  Model Delta
\item  Model Vega
\end{itemize}

Here, an extra cleanse is made. Deltas smaller than 5\% and greater than 95\% in abs value ***LOOK FOR A BETTER APPROACH*** are committed due to potential issues in pricing calculations.


At this point, it would be possible to do a comparison with the Vendor Database

****INCLUDE COMPARISSON ***


\paragraph{Generation of volatility surface}

We implemented the Gaussian Kernel Smoothing


\paragraph{Tests}

Comparison tests were performed between the volatility surfaces obtained following these steps, and a volatility surface provided by a vendor:


\begin{center}
\begin{figure}[h!]
\includegraphics[width=7in]{ALTvsVendor_SPX_04042019}
\caption{Differences in volatilities}
\end{figure}
\end{center}




We see peaks of around 3\% differences for short maturities (1M 3M) and an average of a +-2\% relative difference between both.



\section{Production Work flow}
\begin{itemize}
\item Productionizing the workflow
\item What technology do we use? 
\item Programming language?
\item Platform: Azure + Spark
\item Can we compute IVS on the fly? (real-time)
\end{itemize}

\end{document}